% Responsable: Naranjo Alfonzo Anthony Alexander
% FORM-PC-01 · Ronda 1: caracterización del modelo iterativo.
\subsection{Modelo iterativo}
\label{sec:modelo-iterativo}

El modelo iterativo plantea que los requisitos no tienen que quedar completamente
resueltos en un único intento. Se comienza con una comprensión inicial del
problema y se realizan ciclos de obtención, documentación, negociación y
validación para corregir ambigüedades y ampliar la especificación. Cada ciclo
aprovecha lo aprendido en el anterior; por eso es útil cuando existen
necesidades conocidas, pero todavía quedan detalles por aclarar
\cite{sommerville2016,pohl2025}.

A diferencia de una secuencia estrictamente lineal, una observación del cliente
puede llevar a revisar requisitos ya documentados. El producto de cada ciclo
puede ser una versión mejorada de la Especificación de Requisitos de Software
(ERS), un prototipo evaluado o un conjunto de requisitos acordados. Iterar
sobre requisitos no obliga, por sí solo, a entregar software funcionando en
cada ciclo; esa decisión depende de cómo se organice el desarrollo.

\subsubsection{Fases principales}

\begin{enumerate}
  \item \textbf{Comprensión inicial y delimitación.} Identificar el problema,
  los objetivos, los actores, las restricciones y los requisitos conocidos.
  También se seleccionan los aspectos con mayor incertidumbre.
  \emph{Resultado:} alcance preliminar y primera versión de requisitos.

  \item \textbf{Planificación de la iteración.} Elegir qué funcionalidades,
  dudas o riesgos se analizarán en el ciclo y acordar quiénes revisarán los
  resultados. \emph{Resultado:} objetivos y criterios de revisión del ciclo.

  \item \textbf{Obtención y documentación.} Realizar entrevistas,
  observaciones o revisión de prototipos para precisar necesidades; actualizar
  requisitos funcionales, de calidad y restricciones, identificando sus
  fuentes. \emph{Resultado:} ERS o modelos revisados.

  \item \textbf{Negociación y validación.} Contrastar los requisitos con
  las partes interesadas, resolver contradicciones, comprobar que sean claros
  y verificables, y registrar los acuerdos. \emph{Resultado:} requisitos
  aceptados y observaciones pendientes \cite{ireb2026,pohl2025}.

  \item \textbf{Evaluación y siguiente ciclo.} Revisar las observaciones,
  controlar las versiones, actualizar las prioridades y decidir si se
  necesita otra iteración o si ya existe una especificación suficiente para
  continuar. \emph{Resultado:} nueva línea base o cierre documentado.
\end{enumerate}

La validación y la gestión de cambios acompañan a todo el proceso, no se
reservan únicamente para el último paso. Esto mantiene la trazabilidad entre
las necesidades del cliente, las versiones de los requisitos y las decisiones
tomadas \cite{pohl2025}.

\subsubsection{Supuestos del modelo}

\begin{itemize}
  \item Existen partes interesadas disponibles para revisar los resultados
  en momentos acordados y dar respuestas a las dudas importantes.
  \item Los requisitos pueden precisarse progresivamente sin que el costo de
  cada corrección haga inviable el proyecto.
  \item Es posible dividir el problema en conjuntos de requisitos o asuntos
  prioritarios que puedan analizarse y validarse por ciclos.
  \item El equipo dispone de mecanismos para registrar versiones, decisiones,
  dependencias y solicitudes de cambio, evitando perder acuerdos anteriores.
  \item Se puede reservar tiempo para la retroalimentación y el retrabajo
  que surjan al descubrir necesidades no previstas.
\end{itemize}

\subsubsection{Condiciones de aplicación y límites}

\textbf{Es conveniente} cuando el alcance general se conoce, pero aún hay
incertidumbre sobre pantallas, procesos, reglas de negocio o integración con
otros sistemas. También ayuda cuando los usuarios pueden revisar prototipos o
versiones intermedias y es necesario conservar una especificación formal que
mejore en cada ciclo \cite{sommerville2016}.

\textbf{Es menos conveniente} si no existe ningún representante capaz de
validar avances, si las decisiones dependen de autorizaciones demasiado
lentas o si el proyecto no dispone de tiempo para revisar el trabajo. En
sistemas con interfaces críticas o restricciones contractuales, las
iteraciones deben mantener controles de trazabilidad, aceptación y análisis
de impacto; repetir ciclos no elimina esas obligaciones.

\textbf{Diferencia frente al enfoque ágil.} Ambos aceptan revisiones, pero
el iterativo puede centrarse en refinar la ERS y validar prototipos mediante
ciclos planificados. En el enfoque ágil, la elaboración de requisitos suele
integrarse de forma continua con la construcción y las pruebas de incrementos
pequeños del producto \cite{schon2017}.
